% Bagian: counterfactuals
% Bab: introduction
% Subbagian: counterfactuals

\documentclass[../../../include/open-logic-section]{subfiles}

\begin{document}

\olfileid[id]{cnt}{int}{cnt}

\olsection{Kondisional Kontrafaktual}

Salah satu bentuk konstruksi ``jika \dots maka \dots'' yang sangat umum dan
penting dalam bahasa Inggris dibentuk dengan memakai bentuk subjungtif lampau
dari verba Inggris \emph{to be} (``menjadi/berada''): ``seandainya keadaannya \dots maka keadaannya akan
\dots.'' Karena biasanya anteseden kondisional semacam ini salah, yakni
bertentangan dengan fakta, konstruksi itu disebut \emph{kondisional
  kontrafaktual} (dan karena memakai bentuk subjungtif dari \emph{to be}
(``menjadi/berada''), juga
disebut \emph{kondisional subjungtif}). Kondisional tersebut dibedakan dari
kondisional \emph{indikatif}, yang berbentuk ``jika keadaannya \dots maka
keadaannya \dots.'' Kondisional kontrafaktual dan indikatif mempunyai syarat
kebenaran yang berbeda. Perhatikan contoh terkenal dari Adams:
\begin{quote}
  Jika bukan Oswald yang membunuh Kennedy, orang lainlah yang melakukannya.

  Seandainya Oswald tidak membunuh Kennedy, orang lain tentu akan
  melakukannya.
\end{quote}
Kalimat pertama bersifat indikatif, sedangkan kalimat kedua kontrafaktual.
Kalimat pertama jelas benar: kita tahu bahwa Presiden John F. Kennedy dibunuh
oleh \emph{seseorang}; dan jika orang itu bukan Lee Harvey Oswald (berlawanan
dengan Laporan Warren), maka orang lainlah yang membunuh Kennedy. Kalimat
kedua menyatakan sesuatu yang berbeda. Kalimat itu mengklaim bahwa seandainya
Oswald tidak membunuh Kennedy---yakni, seandainya penembakan di Dallas dapat
dihindari atau tidak berhasil---sejarah kemudian akan berlangsung sedemikian
rupa sehingga percobaan pembunuhan lain akan berhasil. Agar kalimat itu benar,
haruslah ada kekuatan-kekuatan besar yang bersekongkol untuk memastikan
kematian JFK (sebagaimana diyakini oleh banyak penganut teori konspirasi JFK).

Masih menjadi perdebatan apakah kondisional \emph{indikatif} ditangkap secara
tepat oleh kondisional material, khususnya apakah paradoks-paradoks kondisional
material dapat ``dijelaskan'' dengan cara yang sesuai dengan pandangan bahwa
kondisional material memberikan syarat kebenaran bagi kondisional indikatif
dalam bahasa Inggris. Sebaliknya, tidak diperdebatkan bahwa kondisional
kontrafaktual tidak dapat dilambangkan secara tepat dengan kondisional
material. Hal ini jelas karena, meskipun anteseden kontrafaktual pada umumnya
salah, tidak semua kontrafaktual dengan anteseden salah bernilai benar---jika
Anda memercayai Laporan Warren dan tidak ada persekongkolan untuk membunuh JFK,
maka kondisional kontrafaktual Adams di atas merupakan salah satu contohnya.

Kondisional kontrafaktual berperan penting dalam penalaran kausal: salah satu
contoh utama penggunaan kontrafaktual ialah untuk mengungkapkan hubungan
sebab-akibat. Sebagai contoh, menggesek sebatang korek api menyebabkannya
menyala, dan hal ini dapat diungkapkan dengan mengatakan ``seandainya korek
api ini digesek, korek ini akan menyala.'' Kondisional material, dan
kondisional indikatif pada umumnya, tidak dapat dipakai untuk mengungkapkan
hal ini: ``korek api digesek $\lif$ korek api menyala'' bernilai benar jika
korek api tidak pernah digesek, terlepas dari apa yang akan terjadi jika
korek itu digesek. Lebih buruk lagi, ``korek api digesek $\lif$ korek api
berubah menjadi seikat bunga'' juga bernilai benar jika korek itu tidak pernah
digesek, padahal korek api tentu tidak akan berubah menjadi seikat bunga jika
digesek.

Masih diperdebatkan apa tepatnya logika yang benar bagi kontrafaktual. Suatu
analisis kontrafaktual yang berpengaruh diberikan oleh Stalnaker dan Lewis.
Menurut mereka, kontrafaktual ``seandainya~$S$ maka~$T$'' bernilai benar jika
dan hanya jika $T$ benar dalam situasi kontrafaktual (``dunia mungkin'') yang
paling dekat dengan keadaan dunia aktual dan tempat~$S$ benar. Analisis ini
disebut analisis ``ontik'' karena mengacu pada suatu ontologi dunia mungkin.
Analisis-analisis lain memakai probabilitas bersyarat atau teori revisi
keyakinan. Ada sangat banyak usulan logika kontrafaktual yang berbeda. Bahkan
tidak ada satu logika kontrafaktual Lewis--Stalnaker tunggal: meskipun
Stalnaker dan Lewis mengajukan penjelasan yang serupa dengan mengacu pada dunia
mungkin yang paling dekat, asumsi-asumsi mereka menghasilkan inferensi valid
yang berbeda.

\end{document}
